\section{Additional Experiments}\label{app:add_exps}
This section complements the empirical results reported in the main text. We provide the details of our experimental setup, and show the acc/loss heatmaps for several configurations.
\subsection{Empirical Details}
\subsubsection{Toy Example}
In \cref{fig:toy_heatmap}, we trained a simple MLP with LoRA layers to verify the results of the analysis in \cref{sec:toy}. Here we provide the empirical details for these experiments.

\paragraph{Model.} We consider a simple MLP given by
\begin{equation*}
f(x) = W_{out} \phi(BA \phi(W_{in} x)),
\end{equation*}
where $W_{in} \in \reals^{n\times d}, W_{out} \in \reals^{1\times n}, A \in \reals^{r\times n}, B \in \reals^{n \times r}$ are the weights, and $\phi$ is the ReLU activation function. Here, we used $d=5$, $n=100$, and $r=4$.

\paragraph{Dataset.} Synthetic dataset generated by $X\sim \normal(0, I_d), Y = \sin(d^{-1}\sum_{i=1}^d X_i)$ with $d=5$. The number of training examples is $N_{train}=1000$, and the number of test examples is $N_{test}=100$.

\paragraph{Training.} We train the model with gradient descent for a range for values of $(\eta_A,\eta_B)$. The weights are initialized as follows: $W_{in} \sim \normal(0,1.), W_{out} \sim \normal(0, 1/n), A \sim \normal(0, 1/n), B \sim \normal(0, 1.)$. Only the weight matrices $A, B$ are trained and $W_{in}, W_{out}$ are fixed to their initial value. 



\subsubsection{GLUE Tasks with GPT2/Roberta}

For our experiments with GPT2/Roberta-base models, finetuned on GLUE tasks, we use the following setup:

\paragraph{Tasks. } MNLI, QQP, SST2, QNLI

\paragraph{Models.} GPT2, Roberta-base

\paragraph{Training Alg.} AdamW with $\beta_1=0.9, \beta_2=0.99, \epsilon = $ 1e-8, linear schedule, no warmup. 

\paragraph{Learning rate grid.} $\eta_A \in \{$4e-3, 2e-3, 1e-3, 5e-4, 2e-4, 1e-4$\}$, $\eta_B \in \{$ 8e-4, 4e-4, 2e-4, 1e-4, 5e-5, 2e-5, 1e-5 $\}$.

\paragraph{Targert Modules for LoRA.} For Roberta-base, we add LoRA layers to `query' and `value' weights. For GPT2, we add LoRA layers to `c\_attn, c\_proj, c\_fc'.

\paragraph{Other Hyperparameters.} Sequence length $T=128$, train batch size $bs = 32$, number of train epochs $E=3$ ($E=10$ for SST2), number of random seeds $s=3$.


\paragraph{GPUs. } Nvidia V100, Nvidia A10.



\subsubsection{Llama MNLI}
For our experiments using the Llama-7b model, finetuned on MNLI, we use following setup

\paragraph{Training Alg.} AdamW with $\beta_1 = 0.9$, $\beta_2 = 0.999$, $\epsilon =$ 1e-6, constant schedule.

\paragraph{Learning rate grid.} $\eta_A \in \{$1e-6, 5e-6, 1e-5, 2.5e-5, 5e-5, 1e-4$\}$, $\eta_B \in \{$1e-6, 5e-6, 1e-5, 2.5e-5, 5e-5, 1e-4$\}$, $\eta_B \geq \eta_A$

\paragraph{LoRA Hyperparameters.} LoRA rank $r = 8$, $\alpha = 16$, and dropout $0.1$. LoRA target modules `q\_proj, k\_proj, v\_proj, o\_proj, up\_proj, down\_proj, gate\_proj'.

\paragraph{Other Hyperparameters.} Sequence length $T = 128$, train batch size $bs = 32$, number of train epochs $E = 1$, number of random seeds $s = 2$ for $\eta_A = \eta_B$ and $\eta_A, \eta_B$ near test optimal, $s = 1$ otherwise. Precision FP16. 

\paragraph{GPUs.} Nvidia V100.

\subsubsection{Llama flan-v2}

For our experiments using the Llama-7b model, finetuned on a size 100k random subset flan-v2, we use following setup

\paragraph{Training Alg.} AdamW with $\beta_1 = 0.9$, $\beta_2 = 0.999$, $\epsilon =$ 1e-6, constant schedule.

\paragraph{Learning rate grid.} $\eta_A \in \{$1e-6, 5e-6, 1e-5, 2.5e-5, 5e-5, 1e-4$\}$, $\eta_B \in \{$1e-6, 5e-6, 1e-5, 2.5e-5, 5e-5, 1e-4$\}$, $\eta_B \geq \eta_A$

\paragraph{LoRA Hyperparameters.} LoRA rank $r = 64$, $\alpha = 16$, and dropout $0.1$. LoRA target modules `q\_proj, k\_proj, v\_proj, o\_proj, up\_proj, down\_proj, gate\_proj'.

\paragraph{Other Hyperparameters.} Sequence length $T_{\text{source}} = 1536$, $T_{\text{target}} = 512$, train batch size $bs = 16$, number of epochs $E = 1$, number of random seeds $s = 2$ for $\eta_A = \eta_B$ and $\eta_A, \eta_B$ near test optimal, $s = 1$ otherwise. Precision BF16.

\paragraph{MMLU Evaluation.} We evaluate average accuracy on MMLU using 5-shot prompting.

\paragraph{GPUs.} Nvidia A10.

\subsection{Results of Roberta-base Finetuning on all Tasks}
\cref{fig:roberta_main} showed finetuning test accuracy for Roberta-base. To complement these results, we show here the test/train accuracy for all tasks. 

\begin{figure}[h]
    \centering
    \includegraphics[width=0.95\linewidth]{figures/glue_roberta_fp16_main.pdf}
    \caption{GLUE/Roberta-base: same as \cref{fig:roberta_main} with test/train accuracy.}
    \label{fig:enter-label}
\end{figure}

Interestingly, the optimal choice of learning rates for test accuracy differs from that of the train accuracy, although the difference is small. This can be due to mild overfitting occuring during finetuning (the optimal choice of learning rates $(\eta_A, \eta_B)$ for train accuracy probably lead to a some overfitting). 

\subsection{Results of GPT2 Finetuning on all Tasks}
\cref{fig:gpt2_main} showed finetuning results for GPT2 on MNLI and QQP. To complement these results, we show here the test/train accuracy for all tasks. 

\begin{figure}[h]
    \centering
    \includegraphics[width=0.9\linewidth]{figures/glue_gpt2_fp16_all_tasks_halfprecision.pdf}
    \caption{GLUE/GPT2: same setup as \cref{fig:gpt2_main} with additional tasks}
    \label{fig:gpt2_all_tasks_half_precision}
\end{figure}

\newpage
\subsection{GLUE Tasks with Full Precision}


\begin{figure}[h]
    \centering
 \includegraphics[width=0.9\linewidth]{figures/glue_roberta_full_precision.pdf}
    \caption{GLUE/Roberta-base: same as \cref{fig:roberta_main} with full precision training instead of FP16.}
    \label{fig:full_precision_glue_roberta}
\end{figure}



\begin{figure}[h]
    \centering
    \includegraphics[width=0.9\linewidth]{figures/glue_gpt2_all_tasks_fullprecision.pdf}
    \caption{GLUE/GPT2: same setup as \cref{fig:gpt2_all_tasks_half_precision} with full precision training}
    \label{fig:gpt2_all_tasks_full_precision}
\end{figure}

\newpage
\subsection{GLUE Tasks Test/Train Loss}



\begin{figure}[h]
    \centering
    \includegraphics[width=0.9\linewidth]{figures/glue_roberta_half_precision_loss.pdf}
    \caption{GLUE/Roberta-base: same setup as \cref{fig:roberta_main} with $100\times$Test/Train loss instead of accuracy}
    \label{fig:full_precision_glue_roberta_loss}
\end{figure}



\begin{figure}[h]
    \centering
    \includegraphics[width=0.9\linewidth]{figures/glue_gpt2_fp16_all_tasks_loss.pdf}
    \caption{GLUE/GPT2: same setup as \cref{fig:gpt2_all_tasks_half_precision} with $100\times$Test/Train loss instead of accuracy}
    \label{fig:gpt2_halprecision_alltasks_loss}
\end{figure}

\newpage
\subsection{GLUE Tasks with Different LoRA Ranks}

\vspace{-2mm}

\begin{figure}[h!]
    \centering
    \includegraphics[width=0.7\linewidth]{figures/glue_roberta_fullprecision_rank4.pdf}
    \vspace{-2mm}
    \caption{GLUE/Roberta-base: same setup as \cref{fig:roberta_main} with $r=4$}
    \label{fig:roberta_rank4}
\end{figure}

\vspace{-4mm}
\begin{figure}[h!]
    \centering
    \includegraphics[width=0.45\linewidth]{figures/glue_roberta_fullprecision_rank16.pdf}
    \vspace{-3mm}
    \caption{GLUE/Roberta-base: same setup as \cref{fig:roberta_main} with $r=16$}
    \label{fig:roberta_rank16}
\end{figure}

\vspace{-2mm}
\begin{figure}[h!]
    \centering
    \includegraphics[width=0.45\linewidth]{figures/glue_gpt2_fullprecision_rank4.pdf}
    \vspace{-2mm}
    \caption{GLUE/GPT2: same setup as \cref{fig:gpt2_all_tasks_full_precision} with $r=4$}
    \label{fig:glue_rank4}
\end{figure}

\subsection{Experiments with \init{1}}
We also run some experiments using \init{1} as initialization scheme. We noticed that the optimal ratio $\lambda$ is this case is generally smaller than the optimal ratio with \init{2}. \cref{fig:roberta-init1} shows the optimal learning rates $(\eta_A,\eta_B)$ obtained with \init{1} and \init{2}. The optimal ratio $\lambda=\eta_B/\eta_A$ is generally smaller with \init{1}. 

\begin{figure}
    \centering
    \includegraphics[width=0.35\linewidth]{figures/MNLI_init_1.pdf}
    \includegraphics[width=0.35\linewidth]{figures/MNLI_init_2.pdf}
    \caption{Roberta-base with \init{1} and \init{2}, finetuning on MNLI for 10 epochs (similar to \cref{fig:roberta_main} but with more epochs).}
    \label{fig:roberta-init1}
\end{figure}
\newpage
\subsection{Llama Flan-v2 MMLU Acc/Train Loss}
\begin{figure}[!htb]
    \centering
    \begin{subfigure}{0.7\linewidth}
        \includegraphics[width=0.9\linewidth]{figures/llama_flan_v2.pdf}
        \caption{MMLU evaluation accuracy and train loss of Llama-7b trained on flan-v2 100k in the same setting as Figure \ref{fig:llama_test} left panel (using \init{2}). Interestingly, even in one epoch the model can overfit. We were unable to find $\eta_B > \eta_A$ that was optimal for train loss, however it could be the case that the grid was not fine enough or that overfitting does not require much ``feature learning" and $\eta_B / \eta_A \approx 1$ is optimal for minimizing train loss (see the main text for more discussion).\\}
        \label{fig:llama_flan_v2}
    \end{subfigure}
    \begin{subfigure}{0.7\linewidth}
        \includegraphics[width=0.9\linewidth]{figures/llama-7b_flan-v2_init_1.png}
        \caption{MMLU evaluation accuracy and train loss of Llama-7b trained on flan-v2 100k in the same setting as Figure \ref{fig:llama_test} left panel except using \init{1}. Interestingly, the optimal MMLU accuracy is 0.6\% higher than using \init{2} and the optimal ratio $\eta_B / \eta_A$ is twice as large. The training loss is also near optimal only using a large ratio $\eta_B / \eta_A$.}
        \label{fig:llama_flan_v2_init1}
    \end{subfigure}
    \caption{Llama-7b on flan-v2 training with different initializations.}
\end{figure}
\newpage 
\subsection{Llama MNLI Test/Train Loss}

\begin{figure}[h]
    \centering
    \includegraphics[width=0.7\linewidth]{figures/llama_mnli_loss.pdf}
    \caption{Train and test loss of Llama-7b finetuned on MNLI in the same setting as Figure \ref{fig:llama_test} right panel.}
    \label{fig:llama_mnli_loss}
\end{figure}

% \subsection{Finetuning for more epochs}





% \sh{Add Confidence intervals}